\section{The Mystery Spoken Three Times}

\prayeraxis{Announcement and conception}{Luke 1:26--35}{The Father sends; the eternal Son assumes human nature; the Holy Spirit overshadows Mary. The prayer remembers God's initiative, not an angel's independent power.}

\prayeraxis{The handmaid's consent}{Luke 1:38; \textit{Lumen gentium} 56}{Mary receives the word in faith and freely gives the obedience God has graced her to give. Her fiat is real creaturely cooperation, neither divine causality nor passive use.}

\prayeraxis{The Word made flesh}{John 1:1--18, especially 1:14}{The subject of the Incarnation is the eternal Word. He truly becomes human without ceasing to be God, and dwells among us for salvation.}

\subsection{A compressed proclamation of the Incarnation}

The Angelus is named from its Latin opening, but its center is Christ. Gabriel serves God's message; Mary believes and consents; the Word becomes flesh. The Directory on Popular Piety 195 therefore describes the devotion as a recollection of the saving event in which, according to the Father's plan and by the Spirit's power, the Word became flesh in Mary's womb. Its three-part movement is Trinitarian, Christological, and Marian in that order.

The first exchange condenses Luke's annunciation narrative. ``Conceived of the Holy Ghost'' confesses divine initiative and the virginal conception; it does not turn the Spirit into a creaturely father or treat the Incarnation as a biological union between divinity and humanity. The one conceived is the Father's eternal Son, fully divine and fully human.

The second exchange quotes Mary's answer almost verbatim. Her self-description as handmaid places freedom inside discipleship. Vatican II teaches that she was enriched by grace for her vocation and embraced God's saving will wholeheartedly (LG 56). Grace does not abolish her consent; her consent does not place God in her debt.

The third exchange reaches beyond Luke to John's prologue. ``Flesh'' means genuine human existence, not appearance. ``Dwelt among us'' turns contemplation outward: the Son enters human history so that human beings may receive life in him. The Hail Mary after each exchange keeps Gabriel's greeting (Luke 1:28), Elizabeth's blessing (Luke 1:42), and the Church's petition joined to the mystery just proclaimed.

\subsection{The Paschal horizon already inside the prayer}

The collect refuses to isolate Christmas from Easter. The Incarnation made known by an angel leads through Christ's Passion and Cross to the glory of his Resurrection. Paul VI called this Paschal openness one reason the Angelus retained its freshness (\textit{Marialis cultus} 41). The daily pause therefore remembers not a sentimental scene detached from redemption, but the one economy in which the Son takes flesh, gives himself, rises, and draws his people toward glory.

The 1999 \textit{Enchiridion Indulgentiarum} identifies this collect with the Roman Missal's Fourth Sunday of Advent. A liturgical collect incorporated into a pious exercise does not make the whole Angelus a liturgical celebration, and the historical English printed here is not represented as the current vernacular Missal translation.

The sequence also disciplines Christian hope. The prayer asks to be \emph{brought} to resurrection glory. Worthiness for Christ's promises is received through grace and lived in faith, repentance, charity, and perseverance; the vocal act is no automatic mechanism and compels no favor.

\subsection{Intercession under Christ's unique mediation}

``Pray for us'' addresses Mary as a member and mother within Christ's Body, not as a rival redeemer or an independent source of grace. \textit{Lumen gentium} 60--62 and the Catechism 970 make the governing distinction: Christ's unique mediation neither needs completion nor admits a parallel, while created intercession participates wholly by his gift. The Angelus ends by praying to God through Christ. Mary's role is to receive the Word, point to her Son, and intercede within the communion he creates.
